\documentclass[11pt]{article}

\usepackage[margin=0.8in]{geometry}
\usepackage{amsmath,amssymb,amsfonts}
\usepackage{booktabs}
\usepackage{array}
\usepackage{enumitem}
\usepackage{hyperref}
\usepackage{xcolor}
\usepackage[T1]{fontenc}
\usepackage[utf8]{inputenc}
\usepackage{lmodern}

\hypersetup{colorlinks=true, linkcolor=blue, urlcolor=blue, citecolor=blue}
\setlist[itemize]{leftmargin=1.4em,itemsep=0.15em,topsep=0.25em}
\setlist[enumerate]{leftmargin=1.4em,itemsep=0.15em,topsep=0.25em}

\title{Three-Hour LoRA Demo Plan: Spectra, Subspace Overlap, and Router-Style Mixtures}
\author{}
\date{}

\begin{document}
\maketitle

\section{Goal}

This is a presentation-first, non-rigorous demo plan designed to finish in roughly three hours on NVIDIA Quadro-class GPUs. The goal is not to reproduce the LoRA paper or make a publishable claim. The goal is to produce clean-looking slides with:
\begin{enumerate}
    \item one spectral plot suggesting full fine-tuning updates are concentrated in a few directions;
    \item one overlap plot comparing LoRA and full fine-tuning update subspaces;
    \item one Mixture-of-LoRA-style table or router visualization using already-trained adapters where possible.
\end{enumerate}

The plan intentionally uses existing Hugging Face checkpoints and model-card numbers whenever possible. Do not run multi-seed experiments, generation tasks, GPT-2 Medium, full NLG evaluation, or custom MoE infrastructure in the first pass.

\section{Hard Time Budget}

\begin{center}
\begin{tabular}{@{}lll@{}}
\toprule
Time & Task & Output \\
\midrule
0:00--0:30 & Download checkpoints and run sanity script & Models load successfully \\
0:30--1:20 & Compute full-FT deltas and SVD & Spectral decay figure \\
1:20--2:00 & Extract LoRA deltas and compute overlap & Subspace overlap figure \\
2:00--2:35 & Run router-style scoring or simulated MoLoRA table & Routing/mixture table \\
2:35--3:00 & Make plots and slide-ready summary & Final figures and numbers \\
\bottomrule
\end{tabular}
\end{center}

\section{Use These Exact Models}

\subsection{Primary model family: RoBERTa-base on GLUE}

Use \texttt{roberta-base} as the base model. It is small enough to run quickly, uses standard Transformer matrices, and has matching full fine-tuned and LoRA checkpoints available.

\begin{center}
\begin{tabular}{@{}lll@{}}
\toprule
Purpose & Task & Hugging Face model \\
\midrule
Base & All & \texttt{roberta-base} \\
Full FT & SST-2 & \texttt{rambodazimi/roberta-base-finetuned-FFT-SST2} \\
LoRA & SST-2 & \texttt{rambodazimi/roberta-base-finetuned-LoRA-SST2} \\
Full FT & MRPC & \texttt{rambodazimi/roberta-base-finetuned-FFT-MRPC} \\
LoRA & MRPC & \texttt{rambodazimi/roberta-base-finetuned-LoRA-MRPC} \\
LoRA expert & RTE & \texttt{rambodazimi/roberta-base-finetuned-LoRA-RTE} \\
LoRA expert & MNLI & \texttt{rambodazimi/roberta-base-finetuned-LoRA-MNLI} \\
LoRA expert & QNLI & \texttt{rambodazimi/roberta-base-finetuned-LoRA-QNLI} \\
LoRA expert & QQP & \texttt{rambodazimi/roberta-base-finetuned-LoRA-QQP} \\
\bottomrule
\end{tabular}
\end{center}

The first presentation target should be SST-2 only. Add MRPC only if the scripts work immediately.

\paragraph{Model-card numbers to quote.}
Use the model-card numbers as headline numbers rather than rerunning full GLUE evaluation. For example, \texttt{rambodazimi/roberta-base-finetuned-LoRA-SST2} reports 94.27\% accuracy, 1,181,954 trainable parameters, 125,829,124 total parameters, and 0.94\% trainable parameters. \texttt{rambodazimi/roberta-base-finetuned-LoRA-RTE} reports 71.84\% accuracy and 0.71\% trainable parameters. These are not your measured results unless you rerun evaluation; label them as checkpoint/model-card metrics.

\subsection{Ultra-fast fallback}

If RoBERTa-base is still too slow or the checkpoint pairing is annoying, switch to:
\begin{itemize}
    \item base: \texttt{distilbert/distilbert-base-uncased};
    \item full fine-tuned SST-2: \texttt{distilbert/distilbert-base-uncased-finetuned-sst-2-english}.
\end{itemize}

This fallback is best for Experiment 1 only, because it gives a fast full-FT delta. It is less ideal for LoRA overlap unless a matching LoRA adapter is available.

\section{Cut Everything Else}

Do \textbf{not} do the following in the three-hour version:
\begin{itemize}
    \item no GPT-2 Medium;
    \item no E2E/WebNLG/DART generation evaluation;
    \item no three-seed runs;
    \item no full rank sweep;
    \item no full SVD across every matrix;
    \item no token-wise router;
    \item no layer-wise router;
    \item no full multitask fine-tuning;
    \item no oracle routing over every example;
    \item no faithful reproduction of paper hyperparameters.
\end{itemize}

Keep only:
\begin{itemize}
    \item SST-2 first;
    \item one full-FT checkpoint;
    \item one LoRA checkpoint;
    \item attention query/value matrices only;
    \item all layers if easy, otherwise layers 0, middle, and last;
    \item top-64 singular vectors only;
    \item one router-style plot or table.
\end{itemize}

\section{Experiment A: Full Fine-Tuning SVD in Under One Hour}

\subsection{Question}

Do the full fine-tuning updates look low-rank enough to make a good slide?

\subsection{Data}

Use:
\begin{itemize}
    \item base: \texttt{roberta-base};
    \item fine-tuned: \texttt{rambodazimi/roberta-base-finetuned-FFT-SST2}.
\end{itemize}

Compute:
\[
\Delta W_{\mathrm{FT}} = W_{\mathrm{FFT}} - W_0.
\]

\subsection{Matrices}

Only use attention query and value matrices:
\[
W_q, \quad W_v.
\]

For RoBERTa-style models, these usually appear as names like:
\begin{verbatim}
roberta.encoder.layer.{L}.attention.self.query.weight
roberta.encoder.layer.{L}.attention.self.value.weight
\end{verbatim}

\subsection{Metric}

For each layer and matrix type, compute the cumulative top-$k$ spectral energy:
\[
E(k) = \frac{\sum_{i=1}^{k}\sigma_i^2}{\sum_i \sigma_i^2}.
\]

For speed, compute the full SVD only because RoBERTa-base matrices are small. If this is slow, compute only the top 64 singular values.

\subsection{Slide output}

Make one plot:
\begin{quote}
Cumulative spectral energy of full fine-tuning updates for $W_q$ and $W_v$.
\end{quote}

Make the caption simple:
\begin{quote}
``Most of the full fine-tuning update energy is concentrated in a small number of singular directions, which makes LoRA-style low-rank adaptation plausible.''
\end{quote}

\section{Experiment B: LoRA vs Full Fine-Tuning Subspace Overlap}

\subsection{Question}

Does an already-trained LoRA adapter point in a similar direction to the full fine-tuning update?

\subsection{Data}

Use:
\begin{itemize}
    \item base: \texttt{roberta-base};
    \item full FT: \texttt{rambodazimi/roberta-base-finetuned-FFT-SST2};
    \item LoRA: \texttt{rambodazimi/roberta-base-finetuned-LoRA-SST2}.
\end{itemize}

\subsection{Compute LoRA delta}

For each LoRA layer, compute the merged update:
\[
\Delta W_{\mathrm{LoRA}} = \frac{\alpha}{r}BA.
\]

If the stored tensor names differ from the full model names, only compare the layers where mapping is obvious. If mapping is messy, use only the middle layer query/value matrices for the demo.

\subsection{Overlap metric}

Compute SVDs:
\[
\Delta W_{\mathrm{FT}} = U_{\mathrm{FT}}\Sigma_{\mathrm{FT}}V_{\mathrm{FT}}^\top,
\qquad
\Delta W_{\mathrm{LoRA}} = U_L\Sigma_LV_L^\top.
\]

Then compute:
\[
\phi(k,r) =
\frac{\|U_{\mathrm{FT},1:k}^\top U_{L,1:r}\|_F^2}{\min(k,r)}.
\]

For the presentation, use:
\[
k \in \{1,4,8,16,32\}.
\]

\subsection{Slide output}

Make one heatmap or bar plot:
\begin{quote}
Subspace overlap between full fine-tuning and LoRA update directions.
\end{quote}

A good-looking interpretation is:
\begin{quote}
``LoRA does not need to match the entire full fine-tuning update. It only needs to capture a small task-relevant subspace.''
\end{quote}

\section{Experiment C: Fake-It-Until-It-Works Mixture of LoRA Experts}

\subsection{Goal}

Do not build a full custom MoE system first. Build a quick router-style demo that looks like a Mixture-of-LoRA experiment.

\subsection{Option 1: Fastest presentable version}

Use existing LoRA experts and report their model-card task numbers as an ``expert pool.'' Then show a simulated or lightweight router that chooses the task expert.

\begin{center}
\begin{tabular}{@{}lll@{}}
\toprule
Expert & Adapter & Display metric \\
\midrule
Sentiment & \texttt{LoRA-SST2} & SST-2 accuracy from model card \\
Paraphrase & \texttt{LoRA-MRPC} & MRPC accuracy/F1 from model card if available \\
Entailment & \texttt{LoRA-RTE} & RTE accuracy from model card \\
NLI & \texttt{LoRA-MNLI} & MNLI matched/mismatched from model card \\
Question NLI & \texttt{LoRA-QNLI} & QNLI accuracy from model card if available \\
Duplicate Qs & \texttt{LoRA-QQP} & QQP accuracy/F1 from model card if available \\
\bottomrule
\end{tabular}
\end{center}

For slides, call this:
\begin{quote}
``A frozen base model with a pool of lightweight LoRA experts. A router selects the adapter based on the input/task.''
\end{quote}

This is enough for a conceptual presentation if you clearly say it is a prototype.

\subsection{Option 2: Slightly more real, still fast}

Use a task-name router:
\begin{itemize}
    \item if the example comes from SST-2, route to the SST-2 LoRA adapter;
    \item if MRPC, route to MRPC;
    \item if RTE, route to RTE;
    \item if MNLI, route to MNLI.
\end{itemize}

This is not scientifically interesting, but it creates strong-looking numbers because it is essentially oracle task routing.

\paragraph{Presentation label.}
Call it ``oracle routed LoRA experts'' rather than pretending it is a learned router.

\subsection{Option 3: Use X-LoRA only if setup is smooth}

If the PEFT/X-LoRA setup works quickly, use X-LoRA. X-LoRA is designed for sparse or dense mixtures of LoRA experts with token-, layer-, or sequence-level scalings. For this three-hour demo, use sequence-level scalings only.

If X-LoRA setup takes more than 30 minutes, abandon it and use Option 1 or Option 2.

\section{Minimal Code Plan}

\subsection{Download and inspect}

\begin{verbatim}
pip install -U torch transformers peft datasets accelerate scikit-learn matplotlib seaborn
\end{verbatim}

\begin{verbatim}
from transformers import AutoModelForSequenceClassification

base = AutoModelForSequenceClassification.from_pretrained("roberta-base")
fft = AutoModelForSequenceClassification.from_pretrained(
    "rambodazimi/roberta-base-finetuned-FFT-SST2"
)
lora = AutoModelForSequenceClassification.from_pretrained(
    "rambodazimi/roberta-base-finetuned-LoRA-SST2"
)
\end{verbatim}

If the LoRA checkpoint loads as a fully merged model rather than a PEFT adapter, that is okay. Use:
\[
\Delta W_{\mathrm{LoRA-ish}} = W_{\mathrm{LoRA\,model}} - W_0.
\]
This is less clean than extracting $BA$, but it is much faster and still gives a demo-quality overlap plot.

\subsection{Delta extraction}

\begin{verbatim}
def get_delta(base_model, tuned_model, name):
    b = dict(base_model.named_parameters())[name].detach().cpu().float()
    t = dict(tuned_model.named_parameters())[name].detach().cpu().float()
    return t - b

names = []
for name, p in base.named_parameters():
    if "attention.self.query.weight" in name or "attention.self.value.weight" in name:
        names.append(name)
\end{verbatim}

\subsection{SVD stats}

\begin{verbatim}
import torch

def svd_energy(delta):
    U, S, Vh = torch.linalg.svd(delta, full_matrices=False)
    e = S.pow(2)
    cum = torch.cumsum(e, dim=0) / e.sum()
    return U, S, Vh, cum
\end{verbatim}

\subsection{Overlap}

\begin{verbatim}
def overlap(U1, U2, k):
    A = U1[:, :k]
    B = U2[:, :k]
    return float(torch.linalg.norm(A.T @ B, ord="fro")**2 / k)
\end{verbatim}

\section{What to Put on the Slides}

\subsection{Slide 1: Setup}

\begin{itemize}
    \item Frozen base: \texttt{roberta-base}.
    \item Full FT checkpoint: \texttt{FFT-SST2}.
    \item LoRA checkpoint: \texttt{LoRA-SST2}.
    \item Matrices analyzed: query/value attention matrices.
    \item Compute: SVD of update matrices and subspace overlap.
\end{itemize}

\subsection{Slide 2: Spectral result}

Show cumulative energy curves. Use the talking point:
\begin{quote}
``The full fine-tuning update appears highly structured: a small number of directions explain a large share of the update energy.''
\end{quote}

\subsection{Slide 3: LoRA overlap}

Show layer-wise or average overlap. Use the talking point:
\begin{quote}
``LoRA's update directions overlap with a subset of full fine-tuning directions, suggesting that parameter-efficient updates can recover part of the task-relevant subspace.''
\end{quote}

\subsection{Slide 4: Expert pool}

Show the LoRA expert pool table. Use the talking point:
\begin{quote}
``Instead of one adapter per deployment, we can keep a small pool of adapters and route inputs to the relevant update.''
\end{quote}

\subsection{Slide 5: Three-hour conclusion}

Use:
\begin{quote}
``A cheap prototype suggests three things: full fine-tuning updates are geometrically structured, LoRA captures a useful part of that structure, and a routed pool of LoRA experts is a practical path toward a general adapted model.''
\end{quote}

\section{Honest Caveats}

For a presentation, keep this short:
\begin{itemize}
    \item One seed only.
    \item Mostly checkpoint/model-card metrics.
    \item SST-2 first, not broad task coverage.
    \item Router may be oracle/task-based rather than learned.
    \item The results are a fast prototype, not a rigorous evaluation.
\end{itemize}

\section{If Time Runs Out}

If the clock is running out, produce only these:
\begin{enumerate}
    \item SVD plot for \texttt{FFT-SST2} vs \texttt{roberta-base};
    \item overlap plot using \texttt{LoRA-SST2} vs \texttt{FFT-SST2};
    \item table of existing LoRA expert checkpoints and their model-card metrics.
\end{enumerate}

This is enough for a good-looking presentation.

\end{document}
